\documentclass[11pt]{article}
\usepackage[a4paper,margin=1in]{geometry}
\usepackage{fontspec}
\usepackage[boldfont=false]{xeCJK}
\setCJKmainfont{FandolSong-Regular.otf}
\usepackage{amsmath}
\usepackage{amssymb}
\usepackage{array}
\usepackage{booktabs}
\usepackage{graphicx}
\usepackage{adjustbox}
\usepackage{enumitem}
\setlist[itemize]{topsep=2pt,partopsep=0pt,parsep=0pt,itemsep=2pt,leftmargin=1.4em}
\setlist[enumerate]{topsep=2pt,partopsep=0pt,parsep=0pt,itemsep=2pt,leftmargin=1.6em}
\usepackage{parskip}
\usepackage{fvextra}
\fvset{breaklines=true,breakanywhere=true,fontsize=\small}
\setlength{\parindent}{0pt}

\begin{document}

\begin{center}
  {\LARGE\bfseries Inheritance}\\[0.35em]
  {\large A Level Biology}
\end{center}
\vspace{0.6em}

\section*{Meiosis and how it creates variation}

A \textbf{diploid} 二倍体 cell (2n) has two full sets of \textbf{chromosomes} 染色体 — one set from each parent. A \textbf{haploid} 单倍体 cell (n) has just one set.

Chromosomes come in \textbf{homologous chromosomes} 同源染色体 pairs: the two chromosomes in a pair are the same length and carry the same \textbf{genes} 基因 (though they may carry different versions of them).

\textbf{Gametes} 配子 (sex cells) must be haploid. If they were diploid, the chromosome number would double at every generation. So gametes are made by a special division, \textbf{meiosis} 减数分裂, which halves the chromosome number.

Meiosis has two divisions, one after the other, so there are eight named stages: prophase I, metaphase I, anaphase I and telophase I, then prophase II, metaphase II, anaphase II and telophase II. As in mitosis, the \textbf{nuclear envelope} 核膜 breaks down, a \textbf{spindle} 纺锤体 forms, and the chromosomes are moved by the spindle.

Meiosis makes the gametes \textbf{genetically different} from each other in two ways:

\begin{itemize}
\item \textbf{crossing over} 交叉互换 — homologous chromosomes swap matching pieces, mixing the alleles.

\item \textbf{independent assortment} 自由组合 — the pairs line up in a random order, so each gamete gets a random mix of the parent's chromosomes.

\end{itemize}

Then at \textbf{fertilisation} 受精, any gamete can fuse with any other. This random fusion makes every new individual genetically different.

\section*{The language of genetics}

You must know these terms exactly:

\begin{itemize}
\item a \textbf{gene} is a length of DNA that codes for a \textbf{protein} 蛋白质. Its position on a chromosome is its \textbf{locus} 基因座.

\item an \textbf{allele} 等位基因 is one version of a gene.

\item a \textbf{dominant} 显性 allele shows its effect even when only one copy is present; a \textbf{recessive} 隐性 allele only shows when two copies are present.

\item \textbf{codominant} 共显性 alleles both show their effect together.

\item the \textbf{genotype} 基因型 is the alleles an organism has; the \textbf{phenotype} 表现型 is the features you can see.

\item \textbf{homozygous} 纯合 means the two alleles are the same; \textbf{heterozygous} 杂合 means they are different.

\item a \textbf{test cross} 测交 crosses an organism with the recessive homozygote to find its unknown genotype.

\end{itemize}

\section*{Genetic diagrams and crosses}

You show a cross with a genetic diagram, often using a \textbf{Punnett square} 庞纳特方格 — a grid that shows all the ways the gametes can combine.

\begin{itemize}
\item a \textbf{monohybrid cross} 单基因杂交 follows one gene; a \textbf{dihybrid cross} 双基因杂交 follows two genes at once.

\item some genes have \textbf{multiple alleles} 复等位基因 (more than two versions in the population), such as the alleles for human blood groups.

\item in \textbf{sex linkage} 伴性遗传, the gene is on the X chromosome, so the result is different for males and females.

\item in \textbf{linkage} 连锁, genes on the same \textbf{autosome} 常染色体 (a non-sex chromosome) tend to be inherited together.

\item in \textbf{epistasis} 上位性, one gene affects how another gene is shown.

\end{itemize}

\section*{The chi-squared test}

The \textbf{chi-squared test} 卡方检验 compares the results you actually counted with the results you expected from the genetic diagram. It tells you whether the difference is small enough to be due to chance, or large enough to mean something else is going on.

\section*{From genes to proteins to phenotype}

A gene codes for a protein, and that protein affects the phenotype. If the gene is faulty, the protein is faulty, and the phenotype changes:

\begin{center}
\adjustbox{max width=\linewidth}{%
\begin{tabular}{lll}
\toprule
\textbf{Gene} & \textbf{Protein} & \textbf{Effect of a faulty allele} \\
\midrule
TYR & the \textbf{enzyme} 酶 tyrosinase & \textbf{albinism} 白化病 (no pigment made) \\
HBB & \textbf{haemoglobin} 血红蛋白 & \textbf{sickle cell anaemia} 镰状细胞贫血 \\
F8 & factor VIII (helps blood clot) & \textbf{haemophilia} 血友病 \\
HTT & huntingtin & \textbf{Huntington's disease} 亨廷顿病 \\
\bottomrule
\end{tabular}}
\end{center}

\subsection*{Gibberellin and stem height}

In pea plants, the dominant allele \emph{Le} codes for a working enzyme that makes \textbf{gibberellin} 赤霉素, so the plant grows tall by stem \textbf{elongation} 伸长. The recessive allele \emph{le} codes for a broken enzyme, so little gibberellin is made and the plant is short.

\section*{Gene control}

Not all genes are switched on all the time. Cells control which proteins they make.

\begin{itemize}
\item \textbf{structural genes} 结构基因 code for useful proteins such as enzymes; \textbf{regulatory genes} 调节基因 control whether other genes are switched on.

\item an \textbf{inducible enzyme} 可诱导酶 is made only when it is needed; a \textbf{repressible enzyme} 可抑制酶 is normally made but can be switched off.

\end{itemize}

\subsection*{The lac operon}

In a \textbf{prokaryote} 原核生物 such as a bacterium, a group of genes called the \textbf{lac operon} 乳糖操纵子 controls the digestion of lactose:

\begin{itemize}
\item when there is \textbf{no} lactose, a \textbf{repressor} 阻遏物 protein binds to the operon and blocks \textbf{transcription} 转录, so the lactose-digesting enzymes are not made.

\item when lactose \textbf{is} present, it binds to the repressor and pulls it off. Transcription can now happen, and the enzymes are made. The enzymes are therefore inducible.

\end{itemize}

\subsection*{Control in eukaryotes}

In eukaryotes, \textbf{transcription factors} 转录因子 are proteins that bind to DNA and control \textbf{gene expression} 基因表达, by increasing or decreasing the rate of transcription.

Gibberellin switches genes on in this way: it causes the breakdown of \textbf{DELLA} protein repressors. These DELLA proteins normally block the factors that turn on transcription, so removing them lets those genes be expressed.


\end{document}
